        本章关注“已兑现的财务”而非概念标签。2026Q1数据显示，钨和制冷剂标的普遍已经体现收入和利润弹性；电子特气中，中船特气和华特气体等具备产品平台，但按当前股价推算的远期预期显著高于Q1利润基数。

        \begin{exhibitbox}[2026Q1财务交付和当前市值锚]
        \scriptsize
\resizebox{\exhibitwidth}{!}{
\begin{tabular}{lllllll}
\toprule
\textbf{代码} & \textbf{名称} & \textbf{Q1收入} & \textbf{Q1归母净利} & \textbf{毛利率} & \textbf{ROE} & \textbf{估算市值} \\
\midrule
603379 & 三美股份 & 14.0亿元 & 5.1亿元 & 53.4\textbackslash{}\% & 6.0\textbackslash{}\% & 408.1亿元 \\
002842 & 翔鹭钨业 & 11.5亿元 & 2.5亿元 & 31.0\textbackslash{}\% & 18.8\textbackslash{}\% & 170.7亿元 \\
300037 & 新宙邦 & 33.6亿元 & 4.8亿元 & 24.3\textbackslash{}\% & 4.4\textbackslash{}\% & 662.8亿元 \\
605020 & 永和股份 & 13.2亿元 & 1.8亿元 & 27.2\textbackslash{}\% & 3.1\textbackslash{}\% & 210.7亿元 \\
002378 & 章源钨业 & 26.3亿元 & 3.8亿元 & 26.9\textbackslash{}\% & 15.1\textbackslash{}\% & 435.4亿元 \\
600160 & 巨化股份 & 60.2亿元 & 11.7亿元 & 34.4\textbackslash{}\% & 5.6\textbackslash{}\% & 1583.1亿元 \\
002407 & 多氟多 & 32.2亿元 & 3.8亿元 & 27.6\textbackslash{}\% & 4.4\textbackslash{}\% & 645.9亿元 \\
600378 & 昊华科技 & 42.3亿元 & 3.1亿元 & 24.0\textbackslash{}\% & 1.7\textbackslash{}\% & 833.6亿元 \\
300346 & 南大光电 & 6.6亿元 & 1.2亿元 & 41.9\textbackslash{}\% & 3.5\textbackslash{}\% & 607.6亿元 \\
000657 & 中钨高新 & 70.1亿元 & 9.2亿元 & 30.5\textbackslash{}\% & 9.0\textbackslash{}\% & 2459.0亿元 \\
600549 & 厦门钨业 & 157.4亿元 & 11.1亿元 & 21.0\textbackslash{}\% & 6.2\textbackslash{}\% & 2079.9亿元 \\
688549 & 中巨芯-U & 3.7亿元 & 0.1亿元 & 15.7\textbackslash{}\% & 0.2\textbackslash{}\% & 487.1亿元 \\
688146 & 中船特气 & 7.0亿元 & 1.0亿元 & 30.2\textbackslash{}\% & 1.7\textbackslash{}\% & 2041.4亿元 \\
688268 & 华特气体 & 3.8亿元 & 0.3亿元 & 30.2\textbackslash{}\% & 1.7\textbackslash{}\% & 310.1亿元 \\
002971 & 和远气体 & 3.9亿元 & 0.2亿元 & 20.1\textbackslash{}\% & 1.1\textbackslash{}\% & 163.7亿元 \\
603505 & 金石资源 & 9.2亿元 & 0.6亿元 & 14.5\textbackslash{}\% & 3.3\textbackslash{}\% & 210.9亿元 \\
002549 & 凯美特气 & 1.5亿元 & 0.1亿元 & 30.9\textbackslash{}\% & 0.5\textbackslash{}\% & 132.1亿元 \\
688106 & 金宏气体 & 6.6亿元 & 0.0亿元 & 28.1\textbackslash{}\% & 0.1\textbackslash{}\% & 192.2亿元 \\
\bottomrule
\end{tabular}
}
\sourcenote{AStock财务能力抓取；市值为当前价乘以权益/BPS推算股本。}
\end{exhibitbox}

\section{钨链公司：资源弹性高于半导体纯度弹性}
厦门钨业、章源钨业和翔鹭钨业当前首先是钨价重估标的。半导体链条中的高纯钨粉和靶材叙事提高了估值想象，但当期利润主要仍取决于钨精矿、APT、钨粉和硬质合金的价差。中钨高新更偏制造升级和刀具平台，钨价上涨会带来库存和成本扰动，弹性不应简单等同矿端。

\section{氟化工公司：制冷剂是现金流，电子材料是估值期权}
巨化股份、三美股份、永和股份的核心变量是配额、价差和装置调配能力。巨化还通过中巨芯获得电子级HF和含氟电子特气映射，所以估值可给SOTP溢价；三美短期更纯粹地受益制冷剂价差；永和的产业链一体化提供中长期材料弹性。多氟多和新宙邦需要同时处理锂电材料周期、电子化学品认证和有机氟产品结构，不能只按氟化工龙头一把尺子估值。

\section{电子特气公司：直接性和弹性分离}
中船特气是直接性最高的WF6/NF3龙头，但股价已经提前反映远期涨价和国产替代预期。中巨芯-U的直接性强，但盈利基数和订单边界使其更适合事件驱动跟踪。华特气体和南大光电是平台型电子材料公司，品类、客户和认证能力有价值，但并非所有收入都暴露于WF6。金宏气体、凯美特气、和远气体更多是电子气体扩散线，除非看到明确产品和客户，否则应打折处理。
